\documentclass[12pt,spanish,letterpaper,oneside]{article}
\def\documenttitle {Examen Parcial I}
\def\documentsubtitle {Registro documental de la evaluaci\'on}
\def\documentsubject {Examen Parcial I}
\def\documentauthor {Martin Jonathan de la Cruz}
\def\coursename {Antecedentes de los derechos humanos}
\def\coursecode {005}
\def\universityname {Universidad Abierta y a Distancia de M\'exico}
\def\universityfaculty {Licenciatura en Derecho}
\def\universitydepartment {Antecedentes de los derechos humanos}
\def\universitydepartmentimage {departamentos/UnADM}
\def\universitydepartmentimagecfg {height=1.57cm}
\def\universitylocation {Ciudad de M\'exico}
\def\authortable {\begin{tabular}{ll}
Alumno: & Martin Jonathan de la Cruz \\
Matricula: & ES2611202040 \\
Figura docente: & Daniel Alexis Lozano Ortega \\
Reporte interno: & 8 \\
Semana: & 4 \\
Fecha de registro: & 22 de octubre de 2026
\end{tabular}}
\usepackage{iftex}
\ifLuaTeX
\let\pdfpagewidth\pagewidth
\let\pdfpageheight\pageheight
\fi
\input{template}
\usepackage{fontspec}
\setmainfont{arial.ttf}[Path=C:/Windows/Fonts/,BoldFont=arialbd.ttf,ItalicFont=ariali.ttf,BoldItalicFont=arialbi.ttf]
\setsansfont{arial.ttf}[Path=C:/Windows/Fonts/,BoldFont=arialbd.ttf,ItalicFont=ariali.ttf,BoldItalicFont=arialbi.ttf]
\renewcommand{\familydefault}{\sfdefault}
\usepackage{tikz}
\usepackage{eso-pic}
\setcitestyle{authoryear,open={(},close={)}}
\begin{document}
\AddToShipoutPictureBG*{\begin{tikzpicture}[remember picture,overlay]
\node[opacity=0.16,inner sep=0pt] at (current page.center)
{\includegraphics[width=0.72\paperwidth]{img/departamentos/UnADM.pdf}};
\end{tikzpicture}}
\templatePortrait
\templatePagecfg
\templateFinalcfg
\newgeometry{margin=2.5cm}
\fontsize{12}{14.5}\selectfont
\onehalfspacing
\section{Introducci\'on}
El Examen Parcial I aparece como el segundo producto acad\'emico de la
semana 4, junto al cuadro comparativo. Este reporte documenta la evidencia
disponible de su estado; no constituye una resoluci\'on del examen.
\section{Fundamentos, generaciones y protecci\'on jur\'idica de los derechos humanos}
\subsection{Dignidad, reconocimiento y exigibilidad}
La justificaci\'on de un derecho, su reconocimiento normativo y sus
garant\'ias responden a preguntas distintas. El fundamento explica por
qu\'e una persona merece protecci\'on; el reconocimiento permite identificar
facultades y obligaciones; las garant\'ias aportan medios para exigirlas.
La historia de los derechos humanos muestra que la formulaci\'on moderna
de dignidad es resultado de conflictos y transformaciones, no una
realidad siempre reconocida de manera uniforme
\citep[p.~11]{alvarezParcial}.

La perspectiva iusnaturalista sit\'ua el fundamento antes del reconocimiento
estatal; la positivista dirige la atenci\'on a los derechos reconocidos
jur\'idicamente, y la axiol\'ogica a los valores que justifican la protecci\'on.
La lecci\'on de Canal UNED presenta esas perspectivas como aproximaciones
a la fundamentaci\'on \citep{unedParcial}. Ninguna permite deducir por s\'i
sola el resultado de una controversia: el examen jur\'idico requiere normas
aplicables, autoridades obligadas y medios de defensa.

\subsection{Historia y generaciones}
Rey Cantor y coautoras distinguen libertades civiles y pol\'iticas,
derechos econ\'omicos, sociales y culturales y derechos de solidaridad.
Relacionan estos grupos con libertad, igualdad y fraternidad, y examinan
titulares, sujetos obligados y prestaciones. La clasificaci\'on es un
recurso hist\'orico y did\'actico, no una escala de importancia ni una
sucesi\'on que elimine los derechos anteriores
\citep[pp.~72--76]{reyParcial}.

El cap\'itulo sobre primera generaci\'on vincula las libertades civiles
y pol\'iticas con las revoluciones inglesa de 1688--1689, norteamericana
de 1776 y francesa de 1789. Los autores distinguen esas formulaciones
modernas de los antecedentes antiguos y medievales
\citep[p.~83]{reyParcial}. Estudiar antecedentes exige conservar sus
exclusiones y l\'imites, en lugar de atribuirles igualdad universal.

En M\'exico, Brom describe calpulli, autoridades y tributo entre los
mexicas; posteriormente, encomienda, rep\'ublicas de indios y gobierno
virreinal. La presencia de deberes de protecci\'on en un orden desigual
no demuestra su cumplimiento ni equivale al sistema constitucional
actual \citep[caps.~1, 5--6]{bromParcial}. Las transformaciones del
constitucionalismo ampl\'ian tanto la organizaci\'on del poder como los
contenidos individuales y sociales de los derechos
\citep[pp.~1--2]{arteagaParcial}.

\subsection{Obligaciones constitucionales actuales}
El art\'iculo 1 reconoce derechos constitucionales y de tratados de los
que M\'exico es parte. Exige interpretaci\'on conforme y protecci\'on m\'as
amplia, e impone a todas las autoridades promover, respetar, proteger y
garantizar derechos. Universalidad, interdependencia, indivisibilidad y
progresividad impiden convertir la clasificaci\'on generacional en una
jerarqu\'ia. Ante violaciones, el mismo precepto obliga al Estado a
prevenir, investigar, sancionar y reparar en los t\'erminos legales
\citep[art.~1]{cpeumParcial}.

\subsection{Cuestionario: evidencia del intento}
La consulta autenticada del 23 de octubre de 2026 al cuestionario 212011
confirma el nombre \emph{Examen Parcial I}, el indicador \emph{Hecho: Ver}
y el requisito \emph{Por hacer: Hacer intentos: 1}.

\subsection{Condiciones publicadas}
La evaluaci\'on comprende contenidos de las semanas 2, 3 y 4. Contiene
cinco reactivos: cuatro de opci\'on m\'ultiple y uno de verdadero o falso.
Las instrucciones establecen un \'unico intento y un tiempo m\'aximo de
diez minutos, contado desde la selecci\'on de \emph{Comenzar intento}.
Al agotarse ese tiempo se env\'ian autom\'aticamente las respuestas registradas.

La apertura publicada es el 19 de octubre de 2026, 00:00, y el cierre
el 26 de octubre de 2026, 23:55. Se reproducen las fechas visibles del
aula sin atribuir una zona horaria no verificada. Antes de iniciar se
requieren conexi\'on estable, energ\'ia suficiente y un espacio sin
interrupciones. No se inici\'o un intento durante la consulta.

No se dispone en este registro de reactivos, respuestas seleccionadas,
duraci\'on real de un intento, calificaci\'on o retroalimentaci\'on.
El l\'imite de diez minutos es una condici\'on publicada, no una duraci\'on
observada de ejecuci\'on.
La evaluaci\'on se realiza en la plataforma, no mediante un PDF adjunto.
\clearpage
\section{Conclusiones}
La evidencia acredita que el examen est\'a visible y que su requisito de
intento figura pendiente en la captura. Para completar este registro se
requiere la revisi\'on del intento realizada por el estudiante. No se
atribuyen respuestas ni resultados inexistentes.
La distinci\'on entre fundamento, reconocimiento y garant\'ia, junto a la
lectura no jer\'arquica de las generaciones, ofrece una base para estudiar
las semanas 2--4 sin confundir antecedentes con protecci\'on efectiva.
\clearpage
\begin{thebibliography}{6}
\bibitem[\'Alvarez Icaza Longoria(2009)]{alvarezParcial}
\'Alvarez Icaza Longoria, E. (2009). \emph{Los derechos humanos en M\'exico}. Nostra Ediciones.
\bibitem[Arteaga Nava(2014)]{arteagaParcial}
Arteaga Nava, E. (2014). \emph{Derecho constitucional} (4.\textsuperscript{a} ed., edici\'on digital). Oxford University Press.
\bibitem[Brom(2017)]{bromParcial}
Brom, J. (2017). \emph{Esbozo de historia de M\'exico} (1.\textsuperscript{a} ed. digital). Penguin Random House Grupo Editorial.
\bibitem[C\'amara de Diputados(2026)]{cpeumParcial}
C\'amara de Diputados del H. Congreso de la Uni\'on. (2026).
\emph{Constituci\'on Pol\'itica de los Estados Unidos Mexicanos}
(\'ultima reforma DOF 7 de octubre de 2026).
\url{https://www.diputados.gob.mx/LeyesBiblio/pdf/CPEUM.pdf}
\bibitem[Rey Cantor et~al.(2022)]{reyParcial}
Rey Cantor, E., Rey Anaya, G. A., \& Rey Anaya, A. M. (2022).
\emph{Las generaciones de los derechos humanos: Libertad, igualdad y fraternidad}
(10.\textsuperscript{a} ed.). Universidad Libre.
\bibitem[Universidad Nacional de Educaci\'on a Distancia(2024)]{unedParcial}
Universidad Nacional de Educaci\'on a Distancia. (2024, 23 de octubre).
\emph{Tema 11. Los Derechos Humanos Concepto y Fundamentaci\'on}
[Video]. Canal UNED. \url{https://canal.uned.es/video/671b6728e0637d00c70d7b80}
\end{thebibliography}
\end{document}